% =============================================================================
% Ananke-style P-log Recovery for Filesystem Microkernel — Paper Section
% Drop this into your paper's \section or \subsection.
% =============================================================================

\subsection{Persistent Operation Log for Filesystem Recovery}
\label{sec:plog-recovery}

In a disaggregated-memory microkernel, the filesystem server runs as a
user-space process on one machine and may serve clients from other machines via
cross-machine IPC.  If the filesystem's host machine crashes, all in-memory
state---files, directories, and in-flight operations---is lost.  We address
this with a \emph{persistent operation log} (p-log) placed on CXL shared
memory, enabling a surviving machine to transparently recover the filesystem
without any client-visible downtime beyond the in-flight request that was
interrupted by the crash.

\paragraph{Design.}
Each filesystem instance maintains an append-only operation log on a dedicated
1\,MB CXL shared-memory region.  Every mutating operation (file creation,
write) is recorded as a variable-size log entry before the operation returns to
the client.  The log uses a two-phase persistence protocol: (1)~the entry body
is written and flushed to CXL persistence domain via \texttt{clwb}+\texttt{sfence},
then (2)~the header's tail pointer is atomically advanced and flushed.  This
guarantees that on recovery, only fully-written entries are visible---any
partially-written entry beyond \texttt{tail} is simply ignored.

The p-log region is allocated by the kernel during DSM initialization alongside
the existing durable-queue shared-memory regions.  Machine~$N$'s p-log is
identified by a well-known SHM~ID
($\texttt{CLUSTER\_MAX\_MACHINE\_NUM} + N$),
making it discoverable by any machine in the cluster without out-of-band
coordination.

\paragraph{Log format.}
The log begins with a 64-byte cache-line-aligned header containing a magic
number, the current state (\textsc{Active}/\textsc{Crashed}), the tail offset,
and a monotonic sequence counter.  Each entry stores the operation type, a
sequence number, the full pathname, and operation-specific payload:
%
\begin{itemize}
  \item \textsc{Creat}: file path and permission mode.
  \item \textsc{Write}: file path, byte offset, data length, and the written
        data inline (up to 4\,KB per entry, matching the IPC buffer size of the
        durable queue).
\end{itemize}

\paragraph{Instrumentation.}
We instrument the tmpfs server at the VFS operation layer.  On
\texttt{open} with \texttt{O\_CREAT}, a \textsc{Creat} entry is appended
after the file is successfully created.  On \texttt{write}, a \textsc{Write}
entry is appended after the data is committed to the in-memory inode.  To
resolve the file descriptor back to a pathname for write logging, we maintain
a lightweight inode-to-path map that is populated on every \texttt{open}.

\paragraph{Recovery.}
When a machine crashes, a surviving machine starts a recovery filesystem
instance with \texttt{tmpfs.srv --recover $\langle$crashed\_machine\_id$\rangle$}.
The recovery proceeds in four steps:
%
\begin{enumerate}
  \item Initialize an empty tmpfs (root inode, directory structures).
  \item Map the crashed machine's p-log from CXL via \texttt{sys\_mmap\_shm}.
  \item Replay all log entries sequentially: \textsc{Creat} entries invoke the
        internal file-creation path; \textsc{Write} entries resolve the
        pathname to an inode and write the logged data at the recorded offset.
  \item Register the new instance as an IPC server and, optionally, issue
        \texttt{FSM\_REQ\_REPLACE\_FS} to the filesystem multiplexer (FSM) to
        transparently redirect future client connections.
\end{enumerate}

\paragraph{IPC continuity.}
Cross-machine IPC in \sys uses a lock-free durable queue on CXL shared memory
(Section~\ref{sec:appendix-ipc}).  Since the queue itself resides on CXL,
it survives the crash of the consumer machine.  The recovery instance maps the
same queue region and invokes \texttt{dq\_recover\_crash}, which walks the
queue and transitions any node in the \textsc{Doing} state to \textsc{Crash}.
Producers (clients) that were spinning on these nodes observe the
\textsc{Crash} status and can retry.  The recovery instance then starts a new
polling thread on the same queue, seamlessly resuming service for subsequent
client requests.

\paragraph{Crash consistency argument.}
The p-log provides redo-logging semantics.  An entry becomes visible only after
its \texttt{tail} pointer update is persisted; a crash before this point leaves
\texttt{tail} unchanged, so the partial entry is invisible to replay.  Since
each entry is self-contained (includes the full path and data), replay is
idempotent: re-creating an already-existing file returns \texttt{EEXIST} (which
is ignored), and re-writing data at the same offset is a no-op in terms of
final state.  This ensures correct recovery regardless of whether the crash
occurred mid-append or between two operations.
